\section{Problem 4: Stable Matrix}

\subsection{Problem Overview and Objectives}
An $n \times n$ integer matrix $A$ is called \textbf{stable} if every non-boundary (interior) element $A[i][j]$ ($1 \le i, j \le n-2$) equals the sum of its four adjacent interior neighbours:
\[
A[i][j] = A[i-1][j] + A[i+1][j] + A[i][j-1] + A[i][j+1]
\]

If the matrix is not stable, the algorithm must report:
\begin{enumerate}
    \item Total count of violating elements.
    \item Positions $(i, j)$ of all violating interior elements.
    \item Position $(i, j)$ and magnitude of the **maximum absolute difference** $D(i, j)$:
    \[
    D(i, j) = \left| A[i][j] - \left( A[i-1][j] + A[i+1][j] + A[i][j-1] + A[i][j+1] \right) \right|
    \]
\end{enumerate}

\subsection{Specifications}
\begin{itemize}
    \item \textbf{Input}: An $n \times n$ 2D integer matrix \texttt{matrix}.
    \item \textbf{Output}: \texttt{(is\_stable, count, violating\_positions, max\_diff\_position, max\_abs\_diff)}.
    \item \textbf{Constraints}: Matrices with $n < 3$ have no interior elements and are defined as stable by default.
\end{itemize}

\subsection{Requirement of an Additional Matrix}
\begin{quote}
\textbf{Question}: Is an additional matrix required to solve this problem? \\
\textbf{Answer}: \textbf{NO.} An additional matrix is NOT required.
\end{quote}

\textbf{Reasoning}:
\begin{itemize}
    \item Evaluating matrix stability is a read-only process.
    \item Computing neighbour sums for cell $A[i][j]$ reads directly from the original input matrix without modifying any cell values.
    \item Tracking violation counts, coordinates, and maximum absolute differences requires only primitive scalar variables and coordinate tuples.
    \item Therefore, the algorithm operates directly in-place in $O(1)$ auxiliary memory.
\end{itemize}

\subsection{Formal Algorithm Design}

\begin{algorithm}[H]
\caption{Stable Matrix Evaluation and Violation Tracking}
\begin{algorithmic}[1]
\Function{CheckStableMatrix}{$A$}
    \State $n \gets \text{Length}(A)$
    \If{$n < 3$}
        \State \Return $(\text{True}, 0, [], \text{None}, 0)$ \Comment{No interior elements exist}
    \EndIf
    \State $violating\_positions \gets []$
    \State $max\_abs\_diff \gets -1$, $max\_diff\_pos \gets \text{None}$
    \For{$r \gets 1$ \textbf{to} $n - 2$} \Comment{Iterate interior cells only}
        \For{$c \gets 1$ \textbf{to} $n - 2$}
            \State $neighbour\_sum \gets A[r-1][c] + A[r+1][c] + A[r][c-1] + A[r][c+1]$
            \State $val \gets A[r][c]$
            \If{$val \neq neighbour\_sum$}
                \State Append $(r, c)$ to $violating\_positions$
                \State $diff \gets |val - neighbour\_sum|$
                \If{$diff > max\_abs\_diff$}
                    \State $max\_abs\_diff \gets diff$
                    \State $max\_diff\_pos \gets (r, c)$
                \EndIf
            \EndIf
        \EndFor
    \EndFor
    \State $is\_stable \gets (\text{Length}(violating\_positions) == 0)$
    \If{$is\_stable$}
        \State \Return $(\text{True}, 0, [], \text{None}, 0)$
    \Else
        \State \Return $(\text{False}, \text{Length}(violating\_positions), violating\_positions, max\_diff\_pos, max\_abs\_diff)$
    \EndIf
\EndFunction
\end{algorithmic}
\end{algorithm}

\subsection{Mathematical Complexity Analysis}
Let $n$ be the dimension of the $n \times n$ matrix.

\subsubsection{Time Complexity Derivation}
\begin{itemize}
    \item \textbf{Best-Case Time Complexity}: $O(n^2)$.
    Even if a violation is detected at the first interior cell $(1, 1)$, identifying the \emph{global maximum absolute difference} $D(i, j)$ and counting \emph{all} violating elements requires evaluating all $(n-2)^2$ interior cells.
    
    \item \textbf{Worst-Case Time Complexity}: $O(n^2)$.
    There are $(n-2) \times (n-2)$ interior cells. For each cell, performing 4 neighbour additions, 1 subtraction, 1 absolute value calculation, and 1 comparison takes constant $O(1)$ time:
    \[
    T_{\text{worst}}(n) = (n-2)^2 \times O(1) = O(n^2)
    \]
\end{itemize}

\subsubsection{Auxiliary Space Complexity}
\begin{itemize}
    \item \textbf{Auxiliary Space Complexity}: $O(1)$ (excluding output storage for reported violation coordinates). Memory allocation remains constant throughout matrix processing.
\end{itemize}
